% ==========================================
% BAB VII PENUTUP
% ==========================================
\cleardoublepage
\chapter{PENUTUP}
\label{chap:penutup}

\section{Kesimpulan}

Penelitian ini telah mengimplementasikan sistem \textit{Graph Retrieval-Augmented Generation} (GraphRAG) berbasis spesifikasi OpenAPI Kubernetes dan memverifikasi tiga kontribusi teknis utama melalui evaluasi empiris terhadap 97 \textit{fixture} uji, \textit{ablation study}, uji signifikansi statistik menggunakan uji Wilcoxon \textit{signed-rank} dengan \textit{paired bootstrap} dan koreksi Holm-Bonferroni, serta validasi pakar terhadap kualitas jawaban sistem.

\textbf{Pertama}, \textit{pipeline} ekstraksi \textit{knowledge graph} yang deterministik berbasis skema OpenAPI berhasil diimplementasikan. Graf dibangun dari \texttt{swagger. json} versi v1.30 menghasilkan 725 \textit{node} definisi dan 18 jenis \textit{edge} dalam 7 kategori relasi semantik yang khusus dirancang untuk domain Kubernetes. Seluruh \textit{edge} diturunkan dari referensi tipe pada skema sehingga graf bersifat \textit{reproducible} dan dapat diverifikasi secara otomatis, berbeda dengan pendekatan ekstraksi berbasis LLM yang menghasilkan representasi stokastik \parencite{pan2024unifying, wan2025empowering}. \textit{Ablation study} komponen \textit{exact match} (A1) dan \textit{multi-hop traversal} (A2) mengonfirmasi kontribusi tiap-tiap pendekatan, yaitu penghapusan \textit{exact match} menurunkan RetQ sebesar 0,249 poin ($p < 0{,}001$), sedangkan penghapusan \textit{multi-hop} menurunkan RetQ sebesar 0,601 poin ($p < 0{,}001$). Hasil tersebut menunjukkan bahwa \textit{multi-hop traversal} (A2) merupakan komponen paling krusial dan setiap komponen \textit{retrieval} berkontribusi pada hasil jawaban LLM.

\textbf{Kedua}, sistem GraphRAG yang memadukan mekanisme \textit{retrieval} \textit{intent-adaptive depth} dan validasi YAML berbasis \textit{knowledge graph} berhasil dikembangkan. Analisis sensitivitas kedalaman pada $d \in \{1, 2, 3, 4, 5\}$ menunjukkan bahwa tidak ada satu nilai kedalaman tetap yang optimal untuk seluruh tipe \textit{intent}. Ditemukan bahwa tipe \textit{followup} mencapai RetQ optimal pada kedalaman 2 dan menurun tajam pada kedalaman 3, sementara \textit{yaml\_gen} dan \textit{planning} optimal pada kedalaman 3. \textit{Ablation study} A3 (kedalaman tetap $d=2$) mengkonfirmasi bahwa kedalaman seragam yang terlalu dangkal menurunkan RetQ sebesar 0,030 poin dan \textit{Hop Accuracy} sebesar 0,124 poin ($p < 0{,}001$), sedangkan A4 (kedalaman tetap $d=3$) menghasilkan selisih mendekati nol. Perbedaan A3 dan A4 menunjukkan bahwa perilaku adaptif tidak dapat digantikan oleh satu nilai kedalaman tetap. Kontribusi validasi YAML teridentifikasi melalui ablasi A5 yang menunjukkan bahwa penghapusan lapisan validasi Neo4j menurunkan ReaQ secara signifikan ($p = 0{,}007$). Lapisan ketiga ini memverifikasi keberadaan \textit{required fields} langsung dari \textit{knowledge graph} tanpa memerlukan \textit{dry-run} pada kluster aktif.

\textbf{Ketiga}, melalui perbandingan terhadap \textit{Vector RAG} dan \textit{Vanilla LLM}, penelitian ini mengidentifikasi faktor-faktor keunggulan sistem GraphRAG. Pada faktor yang berkaitan dengan \textit{retrieval} dan penalaran berbasis graf, GraphRAG unggul secara signifikan terhadap kedua \textit{baseline} ($p < 0{,}001$, Holm-Bonferroni). Nilai RetQ 0,7259 dibandingkan 0,4255 (\textit{Vector RAG}) dan 0,0241 (\textit{Vanilla LLM}) dengan $\Delta = +0{,}30$ dan 95\% CI $[+0{,}23; +0{,}37]$; ReaQ 0,5554 dibandingkan 0,3307 dan 0,3351; \textit{Path Coverage} 0,8515 dengan $\Delta = +0{,}310$ terhadap \textit{Vector RAG}; serta RGA (\textit{Retrieval-Grounded Answer}) 0,4536 dibandingkan 0,2784 dan 0,0206. Di sisi lain, AnsQ tidak berbeda signifikan antar-sistem ($p_{\text{Holm}} > 0{,}05$) dan hasil \textit{syntactic validity} YAML yang identik antar-model (0,8947) karena kualitas teks jawaban ditentukan oleh kemampuan generatif LLM yang digunakan. Berdasarkan analisis tersebut, keunggulan GraphRAG diidentifikasi terletak pada ketertelusuran (\textit{retrieval-grounded}) dan kelengkapan relasional. Analisis kondisi kritis mengungkap bahwa keunggulan terbesar terjadi pada tipe \textit{intent} relasional (\textit{command} $+0{,}671$, \textit{planning} $+0{,}605$, \textit{followup} $+0{,}574$) dan pada \textit{resource} dengan konektivitas menengah hingga tinggi dalam graf (derajat 3--7, Spearman $\rho = +0{,}467$, $p < 0{,}001$). Validasi pakar memperkuat temuan ini dengan menyatakan \textit{Retrieval Trace} memperoleh tingkat kepercayaan tertinggi (4,50 dari 5), selaras dengan keunggulan RGA yang terukur secara statistik.

\section{Keterbatasan}

Penelitian ini mengidentifikasi tiga keterbatasan utama yang memengaruhi cakupan dan generalisabilitas hasil.

\textbf{Pertama}, \textit{knowledge graph} yang dibangun merupakan \textit{directed graph} yang hanya merepresentasikan relasi struktural-skema dari bagian \texttt{definitions} pada \texttt{swagger.json} Kubernetes v1.30. Keterbatasan ini berdampak pada tiga aspek cakupan: (a) relasi operasional seperti ikatan RBAC antara \textit{ServiceAccount} dan \textit{ClusterRole} tidak dapat dijangkau melalui \textit{traversal} karena \textit{edge} hanya ada dalam arah terbalik (\textit{RoleBinding} $\to$ \textit{ServiceAccount}); (b) operasi API (\texttt{paths}), \textit{Custom Resource Definition}, dan versi Kubernetes selain v1.30 tidak tercakup; serta (c) status \textit{runtime} pada kluster asli Kubernetes tidak terepresentasi dalam graf yang bersifat statis.

\textbf{Kedua}, \textit{retrieval} dibatasi pada maksimum tiga lompatan. Saat kueri membutuhkan informasi di luar jangkauan ini, sistem tidak mengembalikan pesan galat melainkan mengisi kesenjangan konteks menggunakan pengetahuan parametrik LLM (\textit{silent degradation}). Kondisi ini merupakan faktor utama yang menjelaskan mengapa kategori \textit{realworld} memperoleh RetQ terendah (0,39) dan \textit{RetQ-gain} terendah ($+0{,}089$) di antara delapan kategori. Permasalahan ini disebabkan oleh pertanyaan berbasis konteks operasional sering membutuhkan informasi yang melampaui batas tiga lompatan dan di luar cakupan skema statis yang ada pada sistem.

\textbf{Ketiga}, keunggulan GraphRAG tidak menjangkau seluruh dimensi evaluasi. AnsQ tidak menunjukkan perbedaan yang signifikan dari kedua \textit{baseline} dan \textit{syntactic validity} YAML identik (0,8947) di ketiga sistem sehingga keunggulan terbatas pada ketertelusuran dan kelengkapan relasional. Selain itu, nilai \textit{grounding score} 0,7602 menunjukkan bahwa sekitar 24\% istilah Kubernetes dalam jawaban masih berasal dari pengetahuan parametrik LLM.

\section{Saran}

Berdasarkan keterbatasan yang teridentifikasi, beberapa arah pengembangan lanjutan direkomendasikan.

\begin{enumerate}
    \item \textbf{Perluasan Sumber dan Cakupan \textit{Knowledge Graph}} \\
    Pengembangan lanjutan dapat menambahkan relasi operasional seperti ikatan RBAC dan anotasi \textit{best practice} sebagai lapisan kedua di atas graf skema yang sudah ada, serta memperluas ingest ke bagian \texttt{paths} dan \textit{Custom Resource Definition} tanpa mengubah prinsip determinisme graf.

    \item \textbf{Integrasi Sumber Data Dinamis untuk Pertanyaan \textit{Realworld}} \\
    Integrasi dengan \textit{Kubernetes Watch API} atau log operasional dapat meningkatkan kemampuan sistem menjawab pertanyaan \textit{troubleshooting} yang membutuhkan status kluster aktif. Di sisi \textit{retrieval}, strategi kedalaman adaptif untuk kueri lintas-\textit{resource} yang melampaui tiga lompatan dapat mengurangi ketergantungan pada \textit{silent degradation}.

    \item \textbf{Pembaruan Inkremental \textit{Knowledge Graph}} \\
    Kubernetes merilis versi baru secara berkala dengan penambahan atau perubahan skema. Pengembangan \textit{pipeline} pembaruan inkremental yang dapat mendeteksi perubahan antar-versi \texttt{swagger.json} dan memperbarui \textit{knowledge graph} tanpa rekonstruksi penuh akan meningkatkan ketahanan sistem terhadap evolusi API.
\end{enumerate}
